\documentclass[a4paper,12pt]{article}
\usepackage[T2A]{fontenc}
\usepackage[utf8]{inputenc}
\usepackage[russian]{babel}
\usepackage{amsmath,amssymb,mathtools}
\usepackage{helvet}
\renewcommand{\familydefault}{\sfdefault}
\usepackage{icomma}
\usepackage[a4paper,left=22mm,right=22mm,top=23mm,bottom=24mm,headheight=16pt]{geometry}
\usepackage{microtype}
\usepackage{tikz}
\usetikzlibrary{arrows.meta,positioning,shapes.geometric,calc}
\usepackage{booktabs,tabularx}
\usepackage{enumitem}
\usepackage{fancyhdr}
\usepackage{xcolor}
\usepackage{setspace}
\usepackage{needspace}
\usepackage{hyperref}
\definecolor{accent}{HTML}{247C80}
\definecolor{darkaccent}{HTML}{13575C}
\definecolor{textgray}{HTML}{334047}
\definecolor{softgray}{HTML}{617178}
\hypersetup{colorlinks=true,linkcolor=darkaccent}
\setlength{\parindent}{0pt}
\setlength{\parskip}{4pt}
\onehalfspacing
\setlist[enumerate]{leftmargin=*,itemsep=8pt,topsep=8pt}
\pagestyle{fancy}
\fancyhf{}
\fancyhead[L]{\small\color{softgray} ЕГЭ · профильная математика}
\fancyhead[R]{\small\color{softgray} 19.10.2026}
\fancyfoot[L]{\footnotesize\color{softgray} Николь Саркисьянц}
\fancyfoot[R]{\footnotesize\color{softgray}\thepage}
\renewcommand{\headrulewidth}{0.25pt}
\renewcommand{\footrulewidth}{0pt}
\newcommand{\sectionnote}[1]{{\small\color{softgray}#1\par}}
\newcommand{\criterion}[1]{\textcolor{darkaccent}{\bfseries #1}}
\newcommand{\checkitem}[1]{\par\noindent\textcolor{accent}{\(\square\)}\enspace #1\par}
\newcommand{\lessonsection}[1]{\Needspace{5\baselineskip}\section*{\color{darkaccent}#1}\vspace{-5pt}\noindent\textcolor{accent}{\rule{25mm}{0.7pt}}\par\vspace{1pt}}
\newcommand{\Z}{\mathbb Z}
\begin{document}
\thispagestyle{empty}
\begin{titlepage}
\begin{center}
\vspace*{27mm}
{\large\color{darkaccent}\bfseries МАТЕМАТИКА · ЕГЭ}\par
\vspace{27mm}
{\Huge\bfseries Чек-лист}\par
\vspace{9mm}
{\LARGE\color{darkaccent}\bfseries Уравнения и неравенства}\par
\vspace{4mm}
{\Large Тригонометрия, логарифмы и повторение финансов}\par
\vspace{20mm}
\textcolor{accent}{\rule{34mm}{0.7pt}}\par
\vspace{15mm}
{\large Николь Саркисьянц}\par
\vspace{6mm}
{\large 19 октября 2026 года}\par
\vfill
{\small\color{textgray}Лёвин Артём Александрович}\par
{\small\color{textgray}эксклюзивно для Николь Саркисьянц}\par
\vspace{17mm}
\end{center}
\begin{tikzpicture}[remember picture,overlay]
\draw[accent!55,line width=1.1pt] ($(current page.south west)+(0,16mm)$) .. controls ($(current page.south west)+(53mm,29mm)$) and ($(current page.south east)+(-72mm,3mm)$) .. ($(current page.south east)+(0,14mm)$);
\end{tikzpicture}
\end{titlepage}

\lessonsection{01. Дифференцированный кредит: короткое повторение}
\sectionnote{Основной долг погашается равными частями; проценты начисляются на текущий остаток.}
\checkitem{\criterion{Перевожу ставку в доли единицы:} \(p=r/100\), где \(r\)~--- ставка в процентах за период.}
\checkitem{\criterion{Использую формулу общей суммы платежей:} при сумме кредита \(S\) и \(n\) одинаковых периодах}
\[
\boxed{A=S+\frac{Sp(n+1)}2},\qquad
\boxed{p=\frac{2(A-S)}{S(n+1)}}.
\]
\textbf{Пример.} Кредит \(800\) тыс. руб. на два года; общая сумма выплат \(968\) тыс. руб. Найдите годовую ставку.
\[
968=800+\frac{800p\cdot3}{2},\qquad
168=1200p,\qquad p=0{,}14.
\]
\[
\boxed{r=14\%\text{ годовых}}.
\]
\checkitem{\criterion{Проверяю условия применимости:} равномерное уменьшение основного долга и начисление процентов до каждого платежа.}

\lessonsection{02. Показательно-тригонометрическое уравнение}
\sectionnote{Удобная замена превращает составное уравнение в показательное.}
\checkitem{\criterion{Ввожу замену} \(t=\sin^2(3x)\) и сразу фиксирую \(0\le t\le1\).}
\checkitem{\criterion{Применяю свойство степеней:} \(a^{t+m}=a^t\cdot a^m\) при \(a>0\).}

\textbf{Пример.} Решите уравнение
\[
3^{\sin^2(3x)}+3^{2+\sin^2(3x)}
=5^{\sin^2(3x)}+5^{1+\sin^2(3x)}.
\]
Обозначим \(t=\sin^2(3x)\), \(0\le t\le1\). Тогда
\begin{align*}
3^t+9\cdot3^t&=5^t+5\cdot5^t,\\
10\cdot3^t&=6\cdot5^t,\\
\left(\frac35\right)^t&=\frac35.
\end{align*}
Поскольку основание положительно и не равно \(1\), \(t=1\). Обратная замена:
\[
\sin^2(3x)=1\quad\Longleftrightarrow\quad
3x=\frac\pi2+\pi k,\ k\in\Z.
\]
\[
\boxed{x=\frac\pi6+\frac{\pi k}{3},\quad k\in\Z}.
\]
\checkitem{\criterion{Контролирую согласованность:} из \(\sin^2u=1\) следует \(\cos(2u)=-1\), а из \(\sin^2u=\tfrac14\) следует \(\cos(2u)=\tfrac12\). Эти случаи дают разные корни.}

\lessonsection{03. Формулы корней тригонометрических уравнений}
\sectionnote{Эти формулы важно узнавать и применять без построения окружности с нуля.}
\checkitem{\criterion{Помню, что} \(|a|\le1\) необходимо для уравнений \(\sin u=a\) и \(\cos u=a\).}
\renewcommand{\arraystretch}{1.34}
\begin{center}
\small
\begin{tabularx}{\linewidth}{@{}>{\raggedright\arraybackslash}p{0.23\linewidth}X@{}}
\toprule
Уравнение & Полная формула корней, \(k\in\Z\) \\
\midrule
\(\sin u=a\) & \(u=(-1)^k\arcsin a+\pi k\), \(|a|\le1\) \\
\(\cos u=a\) & \(u=\pm\arccos a+2\pi k\), \(|a|\le1\) \\
\(\tg u=a\) & \(u=\arctg a+\pi k\) \\
\(\sin u=0\) & \(u=\pi k\) \\
\(\cos u=0\) & \(u=\dfrac\pi2+\pi k\) \\
\(\sin u=\pm1\) & При \(1\): \(u=\dfrac\pi2+2\pi k\); при \(-1\): \(u=-\dfrac\pi2+2\pi k\) \\
\(\cos u=\pm1\) & При \(1\): \(u=2\pi k\); при \(-1\): \(u=\pi+2\pi k\) \\
\bottomrule
\end{tabularx}
\end{center}
\checkitem{\criterion{Знаю формулу понижения степени:}}
\[
\boxed{\sin^2u=\frac{1-\cos(2u)}2},\qquad
\boxed{\cos^2u=\frac{1+\cos(2u)}2}.
\]
\textbf{Пример.} Из \(\sin^2(3x)=\frac14\) получаем
\[
\frac{1-\cos(6x)}2=\frac14
\ \Longleftrightarrow\ \cos(6x)=\frac12
\ \Longleftrightarrow\ 6x=\pm\frac\pi3+2\pi k.
\]
\[
\boxed{x=\pm\frac\pi{18}+\frac{\pi k}{3},\quad k\in\Z}.
\]
\checkitem{\criterion{Разделяю угол и переменную:} сначала нахожу все значения \(u\), затем делю \emph{каждый член} формулы на коэффициент при \(x\).}

\lessonsection{04. Преобразование выражений: двойной угол и приведение}
\checkitem{\criterion{Распознаю двойной угол:} \(2\sin\alpha\cos\alpha=\sin(2\alpha)\).}
\checkitem{\criterion{Пользуюсь формулами приведения:}}
\[
\sin(180^\circ+\alpha)=-\sin\alpha,\qquad
\sin(180^\circ-\alpha)=\sin\alpha.
\]
\textbf{Пример.} Вычислите выражение
\[
\frac{6\sin95^\circ\cos95^\circ}{\sin170^\circ}
=\frac{3\sin190^\circ}{\sin170^\circ}
=\frac{-3\sin10^\circ}{\sin10^\circ}
=\boxed{-3}.
\]
\checkitem{\criterion{Ищу связанные углы:} \(190^\circ=180^\circ+10^\circ\), \(170^\circ=180^\circ-10^\circ\). Нет необходимости искать только табличные углы.}
\checkitem{\criterion{Пишу знак градуса} у каждого числового угла и проверяю, что знаменатель отличен от нуля.}

\lessonsection{05. Логарифмические неравенства и ОДЗ}
\checkitem{\criterion{Проверяю основание логарифма:} \(a>0\) и \(a\ne1\); аргумент строго положителен.}
\checkitem{\criterion{Помню монотонность:}}
\[
\begin{aligned}
 a>1 &: \quad \log_a U>\log_a V\iff U>V,\\
 0<a<1 &: \quad \log_a U>\log_a V\iff U<V,
\end{aligned}
\quad\text{при }U>0,\ V>0.
\]
\checkitem{\criterion{У вложенного логарифма две ступени ОДЗ:} аргумент внутреннего логарифма положителен, а значение внутреннего логарифма положительно как аргумент внешнего.}

\textbf{Пример.} Решите неравенство
\[
\log_{\sqrt3+1}\!\left(\log_{\sqrt3-1}(2x-3)\right)>0.
\]
Поскольку \(1<\sqrt3<2\), получаем \(0<\sqrt3-1<1\) и \(\sqrt3+1>1\).

\textbf{Шаг 1. ОДЗ.} Необходимо выполнить
\[
2x-3>0,\qquad \log_{\sqrt3-1}(2x-3)>0.
\]
Так как \(0<\sqrt3-1<1\), второе условие означает \(2x-3<1\). Следовательно,
\[
\boxed{\frac32<x<2}.
\]
\textbf{Шаг 2. Решение неравенства.} Основание внешнего логарифма больше \(1\), поэтому
\begin{align*}
\log_{\sqrt3-1}(2x-3)&>1,\\
2x-3&<\sqrt3-1,\\
x&<\frac{\sqrt3+2}{2}.
\end{align*}
При втором переходе знак развёрнут, так как \(0<\sqrt3-1<1\).

\textbf{Шаг 3. Пересечение с ОДЗ.} Из \(1<\sqrt3<2\) следует
\[
\frac32<\frac{\sqrt3+2}{2}<2.
\]
\[
\boxed{x\in\left(\frac32;\,\frac{\sqrt3+2}{2}\right)}.
\]
\checkitem{\criterion{Не штрихую участки вне ОДЗ:} там исходный логарифм не определён.}
\checkitem{\criterion{При умножении или делении на отрицательное число} тоже разворачиваю знак неравенства. При возврате к предыдущему шагу выполняю это правило снова.}

\clearpage
\thispagestyle{fancy}
\begin{center}
{\Large\color{darkaccent}\bfseries Алгоритм: логарифмическое неравенство}\par
\vspace{4mm}
{\color{softgray}От области определения к ответу}\par
\vspace{10mm}
\end{center}
\tikzset{
 aterm/.style={ellipse,draw=darkaccent,line width=.9pt,minimum width=6cm,minimum height=11mm,align=center,font=\small},
 astep/.style={rectangle,rounded corners=2mm,draw=accent,line width=.8pt,text width=10.1cm,minimum height=13mm,align=center,inner sep=4pt,font=\small},
 achoice/.style={diamond,aspect=2.8,draw=darkaccent,line width=.9pt,text width=4.5cm,align=center,inner sep=2pt,font=\small},
 aar/.style={-{Stealth[length=3mm,width=2mm]},line width=.85pt,draw=darkaccent},
}
\begin{center}
\begin{tikzpicture}[node distance=11mm and 22mm]
\node[aterm] (start) {Логарифмическое неравенство};
\node[astep,below=of start] (odz) {Записать ОДЗ: все аргументы \(>0\),\\все основания \(>0\) и \(\ne1\)};
\node[achoice,below=14mm of odz] (choice) {Основание \(a>1\)?};
\node[astep,text width=4.6cm,below left=15mm and 17mm of choice] (yes) {Сохранить знак\\неравенства};
\node[astep,text width=4.6cm,below right=15mm and 17mm of choice] (no) {Развернуть знак\\неравенства};
\coordinate (join) at ($(yes.south)!0.5!(no.south)+(0,-10mm)$);
\node[astep,below=10mm of join] (solve) {Решить полученное неравенство.\\Для следующего логарифма снова проверить основание};
\node[astep,below=of solve] (intersect) {Пересечь решение с ОДЗ\\и проверить граничные точки};
\node[aterm,below=of intersect] (end) {Записать ответ};
\draw[aar] (start)--(odz);
\draw[aar] (odz)--(choice);
\draw[aar] (choice.south) -- ++(0,-5mm) -| node[pos=.78,above,sloped,font=\small]{да} (yes.north);
\draw[aar] (choice.east) -- ++(9mm,0) |- node[pos=.75,above,font=\small]{нет: \(0<a<1\)} (no.north);
\draw[aar] (yes.south) |- (join);
\draw[aar] (no.south) |- (join);
\draw[aar] (join) -- (solve.north);
\draw[aar] (solve)--(intersect);
\draw[aar] (intersect)--(end);
\end{tikzpicture}
\end{center}

\clearpage
\lessonsection{Путь А. Самостоятельная тренировка}
\sectionnote{10 заданий. Выполните без подсказок; затем сверяйте только итоговые ответы.}
\begin{enumerate}
\item Кредит \(600\) тыс. руб. взят на два года с равномерным уменьшением основного долга. Всего выплачено \(726\) тыс. руб. Найдите годовую процентную ставку.
\item Вычислите \(\displaystyle\frac{10\sin80^\circ\cos80^\circ}{\sin20^\circ}\).
\item Вычислите \(\displaystyle\frac{8\sin110^\circ\cos110^\circ}{\sin140^\circ}\).
\item Решите уравнение \(\sin^2(3x)=\dfrac14\).
\item Решите уравнение \(3^{\sin^2(3x)}+3^{2+\sin^2(3x)}=5^{\sin^2(3x)}+5^{1+\sin^2(3x)}\).
\item Решите уравнение \(\displaystyle\left(\frac23\right)^{\sin^2(2x)}=\sqrt{\frac23}\).
\item Решите неравенство \(\log_3(x-1)>1\).
\item Решите неравенство \(\log_{1/2}(4x-3)\ge-1\).
\item Решите неравенство \(\displaystyle\log_{\sqrt3+1}\!\left(\log_{\sqrt3-1}(3x-5)\right)>0\).
\item Решите неравенство \(\displaystyle\log_2\!\left(\log_{1/3}(x-1)\right)\le0\).
\end{enumerate}
\vspace{3mm}
\textbf{Самопроверка:} в № 4--6 укажите период решений; в № 7--10 запишите ОДЗ до преобразований и проверьте включение граничных точек.

\clearpage
\lessonsection{Ответы к тренировочному блоку}
\sectionnote{Параметр \(k\) всюду целочисленный: \(k\in\Z\).}
\renewcommand{\arraystretch}{1.45}
\begin{center}
\begin{tabularx}{\linewidth}{@{}>{\bfseries}p{13mm}X@{}}
\toprule
№ & Ответ \\
\midrule
1 & \(14\%\) годовых. \\
2 & \(5\). \\
3 & \(-4\). \\
4 & \(\displaystyle x=\pm\frac\pi{18}+\frac{\pi k}3\). \\
5 & \(\displaystyle x=\frac\pi6+\frac{\pi k}3\). \\
6 & \(\displaystyle x=\frac\pi8+\frac{\pi k}4\). \\
7 & \(\displaystyle x\in(4;+\infty)\). \\
8 & \(\displaystyle x\in\left(\frac34;\frac54\right]\). \\
9 & \(\displaystyle x\in\left(\frac53;\frac{4+\sqrt3}{3}\right)\). \\
10 & \(\displaystyle x\in\left[\frac43;2\right)\). \\
\bottomrule
\end{tabularx}
\end{center}
\vspace{7mm}
\textcolor{softgray}{\small Особое внимание при повторении: формулы корней тригонометрических уравнений; ограничения при замене; направление знака убывающего логарифма; пересечение решения с ОДЗ.}
\end{document}
